\section{Twierdzenie Petera-Weyla}

\begin{theorem}{Peter-Weyl, 1927}{}
  Niech $G$ będzie grupą zwartą, a $\{e_j\}$ bazą ortonormalną przestrzeni Hilberta $\Hh$. Wtedy współczynniki macierzowe reprezentacji unitarnej $\pi$ grupy $G$ na $\Hh$, czyli odwzorowania
  $$D_{i,j}^{(\pi)}(g)=\langle\pi(g)e_i,e_j\rangle,$$
  generują gęstą podprzestrzeń w zbiorze funkcji ciągłych na $G$.
\end{theorem}


